\clearpage
\section{Results}
This semiconductor laboratory explores the characterization of a \ac{algao} thin film
grown on an r-sapphire substrate. To characterize this system, first the miscut of
the substrate is identified. After that, $\omega$-scans of substrate peaks are conducted
to estimate it's dislocation densities. Subsequent $2\theta\omega$ scans provide
information about the lattice parameters of thin film and substrate. Finally, an
RSM scan visualizes the Ewald sphere.

\subsection{Task 1: Miscut Determination}

\begin{figure}
	\begin{subfigure}
		{0.48\linewidth}
		\centering
		\includegraphics[width=\textwidth]{../plots/task_1_omega.pdf}
		\caption{$\Omega$-scans of the $(024)$ reflection for
		$\varphi=0°,90°,180°,270°$.}
		\label{fig:task1_omega}
	\end{subfigure}
	\begin{subfigure}
		{0.48\linewidth}
		\centering
		\includegraphics[width=\textwidth]{../plots/task_1_phi.pdf}
		\caption{Peak position $\Omega$ as a function of $\varphi$, fitted to determine
		the miscut angle $\gamma$.}
		\label{fig:task1_phi}
	\end{subfigure}
\end{figure}

To measure the miscut angle $\gamma$ of the sapphire substrate, four $\omega$-scans
for different azimutal angles $\phi=0, 90, 180, 270$ were conducted for $\theta=2
6.27°$, measuring the $(024)$ plane. For each scan, the maximum intensity
$I_{\mathrm{max}}$ as well as the corresponding angle $\omega_{\mathrm{max}}$
was extracted, depicted in \cref{fig:task_1_omega}. As discussed as in *section about
miscut*, the $(012)$ reciprocal lattice vector rotates around the $q_{\mathrm{para}}$
axis, resulting in an alternating projection at different $omega$ positions. The data
$\omega(\varphi)$ can be fitted as:
\begin{equation}
	\omega-\theta_{\mathrm{B}}=\arctan[\tan(\gamma)\cos(\varphi+\varphi_{0})]
\end{equation}
It is shown in \cref{fig:task1_phi}. The fit yields an offset angle of $\varphi_{0}
=\qty{-24}{\degree}$ and an miscut angle of $\qty{0.18}{\degree}$.

\clearpage
\subsection{Task 2: Rocking Curves}
$\Omega$-rocking curves were recorded for the symmetric $(006)$ and the asymmetric
$(024)$ reflection to determine the coherence length and the dislocation densities,
see \cref{fig:task2}. The intensity graph of \cref{fig:task2_006} can be reconstructed
using three individual pseudo-voigt functions. A possible explanation of this
complex shape is the formation of three distinct crystallographic domains. Another
interpretation would be that there is the substrate peak, the film peak and one
sub-domain. \Cref{fig:task2_024} displays two well separated peaks. The right
peak corresponds to the substrate's $(024)$ reflection. The left peak could potentially
be an artefact originating from the thin film, as the broad intensity
distribution is captured by the wide-acceptance detector.

The centered peak at \cref{fig:task2_006} yields a center of \qty{20.887}{\degree}
and a FWHM of \qty{0.015}{\degree}. Using *equation to coherence length* yields an
upper boundary for the coherence length of $L_{\parallel}<\qty{902}{\nano \meter}$.
As we fix the dislocation orientation $\mathbf{L}$ parallel to the $(012)$
direction, screw-type dislocations will form parallel to $L$. To find a suitable
burgers vector, we require that $\mathbf{b}$ is a direct lattice vector,
parallel to $\mathbf{L}$. The Ansatz $L_{\parallel}=\frac{0.9 \lambda}{2\Delta\Omega\cdot
\sin(\theta)}$ yields:
\begin{equation}
	u=\lambda'\quad v=2\lambda' \quad w=3 \frac{a^{2}}{c^{2}}\lambda'\simeq 0.4\lambda
	'
\end{equation}
The best approximation yields the $[241]$ direction, with a magnitude of \qty{28}{\angstrom}.
This yields a screw-type dislocation of
$\rho_{\mathrm{screw}}=0.86\cdot 10^{9}m^{-2}$.

The right peak at \cref{fig:task2_024} yields a center of \qty{26.29}{unit} and
a FWHM of \qty{0.01}{\degree}. A suitable burgers vector that is perpendicular to
$L$ is given as $\mathbf{a}_{1}$, since
$\mathbf{a}_{1}\circ (\mathbf{b}_{2}+2\mathbf{b}_{3})=0$ Although the (006)
direction is not perpendicular to (024), one can estimate a edge-type dislocation
of $\rho_{\mathrm{edge}}=73.9\cdot 10^{9}m^{-2}$

\begin{figure}
	\centering
	\begin{subfigure}
		[t]{0.48\linewidth}
		\centering
		\includegraphics[width=\textwidth]{../plots/task_2_006_omega.pdf}
		\caption{$(006)$ reflection}
		\label{fig:task2_006}
	\end{subfigure}
	\begin{subfigure}
		[t]{0.48\linewidth}
		\centering
		\includegraphics[width=\textwidth]{../plots/task_2_024_omega.pdf}
		\caption{$(024)$ reflection}
		\label{fig:task2_024}
	\end{subfigure}
	\caption{$\Omega$-rocking curves of the symmetric $(006)$ and asymmetric
	$(024)$ reflections, each fitted with pseudo-Voigt peaks.}
	\label{fig:task2}
\end{figure}

\subsection{Task 3: Lattice Parameter Determination}

$2\theta\text{-}\Omega$-scans of the $(006)$ and $(036)$ reflections were used
to determine the lattice constants $a$ and $c$, and $\Omega$-rocking curves of both
reflections were compared, see \cref{fig:task3}.

As stated as in *section of coordinate transformation*, the scattering vector length
can be calculated as $q=4\pi /\lambda \cdot \sin(\theta)$. For the $(006)$
reflex, the scattering vector is related to the lattice constant as $\lvert \mathbf{q}
\rvert = 6 \lvert \mathbf{b}_{3}\rvert =6 \cdot 2\pi /c$ This leads to a lattice constant
of $c=1.29929\pm0.00001 nm$. Another small peak is visible at a shifted
$2\theta$. This yields a lattice constant of $c=1.3042\pm0.0005 nm$ that could be
the thin film.

The $(036)$ reflex is related to both $a$ and $c$ lattice constants of the
substrate as
\[
	\lvert \mathbf{q}\rvert =\sqrt{ 3 \cdot \frac{4 \pi}{\sqrt{ 3 }a} + 6 \cdot \frac{2\pi}{c}
	}
\]
Substituting $c$ from above yields an $a=4.75925\pm0.00001 nm$. No thin film $a$ constant
can be found in the spectrum.

\begin{figure}
	\centering
	\begin{subfigure}
		[t]{0.48\linewidth}
		\centering
		\includegraphics[width=\textwidth]{../plots/task_3_006_2to.pdf}
		\caption{$2\theta\text{-}\Omega$-scan of the $(006)$ reflection}
		\label{fig:task3_006_2to}
	\end{subfigure}
	\begin{subfigure}
		[t]{0.48\linewidth}
		\centering
		\includegraphics[width=\textwidth]{../plots/task_3_036_2to.pdf}
		\caption{$2\theta\text{-}\Omega$-scan of the $(036)$ reflection}
		\label{fig:task3_036_2to}
	\end{subfigure}
	\begin{subfigure}
		[t]{0.48\linewidth}
		\centering
		\includegraphics[width=\textwidth]{../plots/task_3_omega.pdf}
		\caption{$\Omega$-rocking curves of the $(006)$ and $(036)$ reflections}
		\label{fig:task3_omega}
	\end{subfigure}
	\caption{$2\theta\text{-}\Omega$-scans of the $(006)$ and $(036)$ reflections
	used to determine the lattice constants $a$ and $c$, and a comparison of the corresponding
	$\Omega$-rocking curves.}
	\label{fig:task3}
\end{figure}

\subsection{Task 4: Reciprocal Space Map}

A frame-based reciprocal space map was recorded around a single reflection, see
\cref{fig:task4_rsm}.

\begin{figure}[H]
	\centering
	\includegraphics[width=\linewidth]{../plots/task_4_rsm.pdf}
	\caption{Reciprocal space map showing the full scanned region and a zoom onto
	the reflection.}
	\label{fig:task4_rsm}
\end{figure}